\section{Deleting certified identity steps before decoder compilation}
\label{sec:identity-collapse}

The six excluded guarded cases suggest a source transformation that
is different from a permutation of all leaves.  Each excluded group
contains the positive branch of an always-zero decrement and its
unreachable loop body.  Removing the whole loop also removes its
feasible zero test, which is an identity operation.  The old and new
source runs then take different numbers of source steps, so the
equal-step leaf-map theorem alone does not apply.

Let $G$ be the old register graph with the certified product invariant,
and choose a set $D$ of leaves to delete.  Require every deleted leaf
with nonempty invariant to be either an unconditional control operation
or a decrement whose tested register is forced to zero by its mask.
Such a deleted step changes no register.  Leaves with empty invariant
may have arbitrary operations, since invariant preservation prevents
their execution.  Require the directed graph of feasible deleted steps
to contain no cycle.  Equivalently, assign deleted feasible leaves a
natural-valued rank that strictly decreases whenever the next leaf is
also deleted.  This requirement rules out an infinite run consisting
solely of steps that the transformation removes.

Starting from any reachable old leaf, repeatedly execute deleted
identity steps until reaching a retained leaf.  Write $\nu(s)$ for
the resulting retained leaf.  It is independent of the particular
register values in the certified source product, because every feasible
deleted decrement has only its zero guard.  The rank proves that this
normalization terminates and leaves all registers unchanged.

\begin{theorem}[Invariant identity-collapse certificate]
Let $H$ be a new register graph and let $f$ bijectively map retained old
leaves to new leaves.  Suppose $f$ sends the normalized old initial leaf to the
new initial leaf, preserves retained register effects and halt tags,
and, for every feasible guarded edge $s\to t$ from a retained old leaf,
the corresponding new edge has target $f(\nu(t))$.  Then $G$ and $H$
halt on exactly the same initial register files in the certified
domain.  The invariant transported to new leaf $f(s)$ is the old
invariant at $s$.
\end{theorem}
\begin{proof}
Normalize the old initial leaf, taking finitely many identity steps.
At a pair of corresponding retained and new leaves the registers agree,
so the same decrement guard is selected.  Their common effect is
followed in the old graph by finitely many deleted identity steps.
The target equation restores the correspondence, and the old invariant
still holds by induction.  Every new nonhalting source step therefore
matches at least one old step and finitely many additional identity
steps.  Halt tags agree at the related leaves.  Concatenating these
finite segments proves one implication of finite halting; deleting the
finite intervening identity segments from an old finite halt proves
the other.  Termination of normalization prevents an infinite deleted
segment from invalidating the correspondence.
\end{proof}

The accompanying \texttt{collapse\_cfg.py} checks this criterion for
leaf-order matching of the retained sites.  Its default deletion set
is
\[
 D=\{6316,6317,6318,6434,6435,6436\},
\]
the two loop triples identified in
Section~\ref{sec:guard-invariant-transport}.  It independently checks
the original product certificate, normalizes deleted feasible paths,
rejects a deleted cycle or a feasible register change, and compares
every feasible retained successor with the new graph.  Break and
continue operands may be retargeted provided their actual leaf targets
pass the check.  The criterion is an ordinary finite proof in ZFC;
it does not itself certify a changed binary machine or its initializer.

The three delivered source experiments pass this independent criterion:
\[
\begin{array}{l|r|r|r}
\text{deleted loop}&\text{retained leaves}&
\text{feasible retained cases}&\text{maximum skipped identity steps}\\\hline
X\text{-transfer}&6457&8384&1\\
T\text{-transfer}&6457&8384&1\\
\text{both transfers}&6454&8383&1
\end{array}
\]
All preserve initial leaf zero.  Their site lists are supplied as
\path{verification/fusion/zero_specialization/}
\texttt{collapseX.sites.json}, \texttt{collapseT.sites.json}, and
\texttt{collapseXT.sites.json}.  Their independent certificates are
\texttt{verification/proof/collapseX-independent.json} and the
corresponding \texttt{T} and \texttt{XT} files.  This establishes a
source-level optimization on the stated invariant.  It does not turn
the costs of the experimentally emitted tables into binary-machine
correctness theorems or supply a sub-$246$ construction.
